% Chapter 4. Add a section for every new experiment and keep the status column of
% tab:experiments current (done / running / planned).
\chapter{Experimental design}
\label{ch:experiments}

This chapter defines the experiments that answer the research questions of \cref{sec:rq} and test
the hypotheses of \cref{sec:hypotheses}. \Cref{tab:experiments} gives an overview, and the
sections below give the design of each experiment: what is varied, what is held fixed, and what
is measured.

\begin{table}[ht]
  \centering
  \caption{Experiments, the questions they address and their status at the time of writing.}
  \label{tab:experiments}
  \small
  \begin{tabularx}{\textwidth}{@{}lLlll@{}}
    \toprule
    & Experiment & RQ & Hypothesis & Status \\
    \midrule
    E1 & LLM designs GQDs as hexagon lists: prompt \(\times\) feedback ablation & RQ1 & H1 & done \\
    E2 & Calibration of the sensing score on 50 dataset GQDs & RQ2, RQ4 & H2, H4 & done \\
    E3 & xTB against DFT on two small models (tetracene, coronene) & RQ2 & H2 & done \\
    E4 & xTB against DFT on 12 GQDs (GPU benchmark) & RQ2 & H2 & done \\
    E5a & Pilot: LLM against random mutation, one run each & RQ3, RQ4 & H3, H4 & done \\
    E5b & Main comparison: three runs per mutator, loop version 2 & RQ3, RQ4 & H3, H4 & done \\
    E6 & Baseline: random GQDs without evolution, same xTB budget & RQ3 & H3 & planned \\
    \bottomrule
  \end{tabularx}
\end{table}

% ==========================================================================================
\section{E1: LLM design with hexagon coordinates}
\label{sec:e1}

\paragraph{Question.} How often does the LLM produce a valid GQD when it writes hexagon positions,
and how much do the closed-shell rule in the prompt and the fix edits in the feedback help?

\paragraph{Requests.} Each request asks for one GQD with a target size drawn at random between 20
and 40 hexagons and one of eight shape descriptions: roughly round, elongated (a ribbon about three
hexagons wide), triangular, rectangular, irregular, mostly armchair edges, mostly zigzag edges, or
L-shaped. The same 80 requests (fixed random seed) are used in every condition. The requests go to
the same LLM deployment, tool and validation code as the agent. When a design fails, the rejection
message is returned and the LLM may retry up to two times.

\paragraph{Conditions.} A \(2 \times 2\) design varies the prompt and the feedback:
\begin{itemize}
  \item prompt: \emph{basic} (coordinates, neighbour rule and size only) or \emph{rule} (basic
        plus the closed-shell rule and the point-symmetry tip of \cref{sec:rule-prompt});
  \item feedback: \emph{plain} (``no Kekul\'e form'') or \emph{fixes} (the imbalance and the fix
        edits of \cref{alg:fixes}).
\end{itemize}
An earlier run of 50 requests with the basic prompt and plain feedback was used to find the main
cause of failure before the ablation was designed.

\paragraph{Measures.} Share of requests valid at the first attempt and after retries, mean number
of attempts, failed attempts by reason, number of distinct valid designs and the most repeated
design, novelty against the 91{,}174-structure dataset, share of point-symmetric designs, mean
deviation from the requested size, predicted gap range, and the number of designs with holes.

% ==========================================================================================
\section{E2: calibration of the sensing score}
\label{sec:e2}

\paragraph{Question.} How do the \ce{NO2} descriptors and the score vary across GQDs, and which
GQDs score best?

\paragraph{Sample.} Fifty GQDs were drawn from the 14{,}674 xTB-labelled structures, five from
each 0.2~eV bin of the dataset gap between 0 and 2.03~eV (random seed 0). They have 20 to 37
hexagons.

\paragraph{Procedure.} Each GQD was evaluated with all five gases (\cref{alg:evaluate}). As a
check of the geometry step, the gap of each optimised GQD was compared with the gap stored in the
dataset.

\paragraph{Measures.} Adsorption energy, charge transfer and chemisorption of \ce{NO2}; adsorption
energies of the interfering gases; selectivity margin; score; and their dependence on the gap. A
small classification study tested whether simple structural features predict chemisorption (area
under the ROC curve).

% ==========================================================================================
\section{E3: xTB against DFT on small models}
\label{sec:e3}

\paragraph{Question.} Where do GFN2-xTB adsorption energies and charges agree with DFT?

\paragraph{Systems.} Tetracene, on which xTB chemisorbs \ce{NO2}, and coronene, on which xTB
physisorbs it. These two cover both binding regimes and are small enough for DFT on the laptop.

\paragraph{Procedure.} The xTB-optimised complexes of all gases at both sites were recomputed with
DFT (\cref{sec:dft-methods}) at the same geometries.

\paragraph{Measures.} Adsorption energy and charge on the gas for every complex, from xTB and from
DFT.

% ==========================================================================================
\section{E4: xTB against DFT on twelve GQDs}
\label{sec:e4}

\paragraph{Question.} Does xTB rank full-size GQDs in the same order as DFT for the quantities
that enter the score, and do the best GQDs bind \ce{O2} in DFT?

\paragraph{Sample.} Ten calibration GQDs spread over the range of scores, and the best GQD of each
pilot run (\code{llm\_s0\_g10c004} and \code{random\_s0\_g02c004}).

\paragraph{Procedure.} 162 DFT single points on one A100 GPU, in three phases
(\cref{sec:dft-methods}).

\paragraph{Measures.} For each complex, the xTB and DFT adsorption energy and charge. For each GQD,
the descriptors and score from both methods. The main outcome is Spearman's rank correlation
\(\rho\) \cite{spearman1904} between xTB and DFT for \(E_{\ce{NO2}}\), \(q_{\ce{NO2}}\) and the
score. We fixed the decision rule before seeing the results:
\begin{itemize}
  \item \(\rho \geq 0.7\): xTB ranking is accepted, and the loop results stand as they are;
  \item \(0.4 \leq \rho < 0.7\): xTB is used for screening, and the top GQDs are re-ranked with
        DFT;
  \item \(\rho < 0.4\): the score must change, for example to a response term based on the HOMO
        energy or a correction learned from DFT.
\end{itemize}
The \ce{O2} adsorption energies give a first check of air stability.

% ==========================================================================================
\section{E5: LLM against random mutation}
\label{sec:e5}

\paragraph{Question.} With the same number of xTB evaluations, does the LLM mutator find better
GQDs than the random mutator?

\paragraph{Fixed across runs.} The starting pool (the 50 calibration GQDs; best score 0.322), the
validation and repair, the gap pre-filter, the number of generations (10) and of xTB-scored
children per generation (up to 10), the selection rule, and the final check. The xTB budget is
therefore at most 100 evaluations per run for both mutators.

\paragraph{Varied.} The mutator, and the random seed. In the pilot (E5a) each mutator ran once
with seed 0 and the first version of the loop. In the main comparison (E5b) each mutator runs three
times with seeds 1, 2 and 3 and version 2 of the loop (\cref{tab:loop-versions}). The seed fixes
all random choices of the loop. It does not fix the LLM, whose replies vary from call to call; for
the LLM runs the seed fixes only the loop around it.

\paragraph{Measures.} For each run: the number of proposals; the share that were valid; the share
that became valid only after repair by code; the number scored by xTB; the share of scored children
that beat their parent; the share that chemisorbed \ce{NO2}; the mean score of the scored children;
the best score; and the number of children above the best starting score. The curve of the best
score against the number of xTB evaluations shows sample efficiency. The number of LLM calls is
reported separately, because LLM calls are not part of the xTB budget.

\paragraph{Statistics.} For each mutator we report the mean and standard deviation over runs of
the best score and of the mean score. As a rough guide, a one-sided Mann-Whitney U test
\cite{mannwhitney1947} compares the scores of all scored children of the two mutators; children of
the same run are not independent, so this test overstates the evidence. With three runs per
mutator, a rank test on the per-run best scores cannot give a one-sided \(p\)-value below 0.05
(there are only 20 ways to split six runs into two groups of three). We therefore judge the main
comparison by the size of the differences and by whether they hold in every pair of runs, and not
by a significance threshold alone.

% ==========================================================================================
\section{E6: baseline without evolution (planned)}
\label{sec:e6}

The original plan compares three methods with the same xTB budget: random generation of new GQDs
without any evolution, evolution with random edits, and evolution with LLM edits. E5 covers the
last two. E6 will generate GQDs with the same random growth procedure that produced the dataset,
apply the same validation and pre-filter, and score about 100 of them with xTB. It shows how much
the evolutionary loop itself adds, independent of the mutator.
\pending{Run E6 after E5b finishes.}
